% Level 1, Week 2 -- answer key
\answersweek{2}{Loop-de-Loop}{Answers for workbook chapter 2 and homework 2.}

\section{Workbook questions}
\begin{answer}[Turn table]
\begin{tabular}{@{}lccc@{}}
\textbf{Shape} & \textbf{Sides} & \textbf{Repeat} & \textbf{Turn} \\
Triangle & 3 & 3 & 120 \\
Square & 4 & 4 & 90 \\
Pentagon & 5 & 5 & 72 \\
Hexagon & 6 & 6 & 60 \\
Octagon & 8 & 8 & 45 \\
\end{tabular}

Rule: repeat = number of sides; turn = 360 ÷ sides. Accept \enquote{( 360 / 5 )} written with the operator block.
\end{answer}

\begin{answer}[Exit ticket]
Pentagon: repeat 5 times, turn 72 degrees.
\end{answer}

\section{Workbook exercises}
\begin{answer}[Exercise 2.1: Walking Cat]
\textbf{Expected behaviour:} after the green flag, the cat walks across the stage, its legs move (it flips between costume1 and costume2 every 0.1 seconds), it turns round at the edge without going upside down, and it keeps walking until the red stop sign is clicked.

\textbf{Question \enquote{what happens without wait?}} The cat zooms and the costumes flicker so fast that the walk looks blurry. The wait slows the loop so we can see each costume.

\textbf{Marking:} \tickbox{} uses \enquote{forever} \quad \tickbox{} next costume inside the loop \quad \tickbox{} bounces at the edge \quad \tickbox{} rotation style left-right (if missing, the cat walks upside down after bouncing: that is a good bug to discuss, not a fail).
\end{answer}

\begin{answer}[Exercise 2.2: Polygon Party]
\textbf{Expected behaviour:} the square script erases the stage and draws a closed square with sides of 100 steps, starting in the middle and turning right (so the square hangs below and to the right of the centre). The student then changes the numbers:
triangle = repeat 3, turn 120; hexagon = repeat 6, turn 60; pentagon = repeat 5, turn 72. Each shape must close.

\textbf{Leo's bug:} repeat 3, turn 60 draws three sides of a hexagon: an open, bowl-like shape that does not close. The bug is the turn: 60 is the turn for 6 sides. The fix is turn 120 (360 ÷ 3), so the three turns add up to 360.

\textbf{What to look for:} students predict before clicking; move and turn are both inside the repeat; \enquote{pen down} is before the loop and \enquote{pen up} after it.
\end{answer}

\begin{answer}[Exercise 2.3: Star Burst]
\textbf{Expected behaviour:} 36 long lines criss-cross to make a round, spiky star, about 150 steps wide, centred near the middle of the stage. The pen colour changes a little after every line, so the star is a rainbow. It ends where it started (36 × 170 = 6120 = 17 whole turns).

\textbf{Open-ended part:} any numbers are fine. Good discoveries to praise: turns close to 180 make spiky stars; turns that divide nicely into 360 make simple shapes; a small move with a small turn makes a circle.
\end{answer}

\section{Homework}
\begin{answer}[Core: Shape Gallery]
Hexagon numbers: repeat \textbf{6}, turn \textbf{60}.

\textbf{Expected behaviour:} three closed shapes in a row near the top of the stage (a triangle on the left, a square in the middle, a hexagon on the right), with no lines joining them. Each shape hangs below and to the right of its start point because the cat turns right.

\textbf{Common slips:} a line joins the shapes (missing \enquote{pen up} before a \enquote{go to}); the hexagon is open (turn is not 60).
\end{answer}

\begin{answer}[Challenge: the circle trick]
It looks like a circle. It is really a 36-sided shape, but the sides are so short (10 steps) and the turns so small (10 degrees) that our eyes see a smooth curve. 36 × 10 = 360, so it goes all the way round once and closes.
A 10-sided shape needs turn 36 (360 ÷ 10).
\end{answer}

\begin{answer}[Creative: loop art]
Any picture that uses at least one loop. Look for loops with something changing inside them (pen colour, pen size, a turn after each shape). A sketch that matches the screen shows good observation.
\end{answer}
